\documentclass{article}
\usepackage[utf8]{inputenc}
\usepackage{ graphicx }
\graphicspath{ {img/} }
\usepackage{ geometry}
\geometry{legalpaper, margin=1.3in}
\usepackage{ bbm }
\usepackage{ braket }
\usepackage{ mathrsfs }
\usepackage{ amssymb }
\usepackage{ mathtools }
\usepackage{ amsfonts }
\usepackage{ caption }
\usepackage{ enumitem } % Lists bullet
\usepackage{ tikz }

\usetikzlibrary{ quantikz }

\newcommand\eqdef{\mathrel{\overset{\makebox[0pt]{\mbox{\normalfont\tiny\sffamily def}}}{=}}}
\newcommand\eqgoal{\mathrel{\overset{\makebox[0pt]{\mbox{\normalfont\tiny\sffamily goal}}}{=}}}

\title{qk4cs}
\author{}
\date{2021/2022}

\begin{document}

\begin{titlepage}
\begin{center}

\textsc{\LARGE Sorbonne Universite} \\[1cm]
\textsc{\LARGE Master 1 - Quantum Information} \\[1cm]
\textsc{\Large 2021/2022} \\[7cm]


{\huge \bfseries Lecture notes} \\[0.5cm]
{\huge \bfseries -} \\[0.5cm]
{\huge \bfseries Quantum Kinematic} \\[4cm]
\vfill
\includegraphics[scale=0.6]{su.png}
\end{center}
\end{titlepage}

\tableofcontents

\newpage

\section{Introduction}
Physical system which has $d \in \mathbbm{N}$ possible distinguishable states.
Its physical state $\ket{\psi} \in \mathscr{H}$, the Hilbert space $\mathbbm{C}^d$.

\begin{equation}
    \ket{\psi} = \begin{bmatrix} \psi_1\\ \psi_2 \\ \vdots \\ \psi_d\end{bmatrix}
    \text{ and } \forall i, \psi_i \in \mathbbm{C}.
\end{equation}


\noindent
The result of the measurement in the computational basis on
$\ket{\psi}$
 is $i \in [1, \cdots, d]$ with probability $|\psi_i|^2$.
 \\ \noindent
And $\sum_{i=1}^d |\psi_i|^2 = \braket{\psi|\psi} = 1$: the state is normalized.

\subsection{Dirac notation}

\begin{itemize}[label=-]

\item Ket:
\begin{equation}
\ket{\psi} = \begin{bmatrix} \psi_1\\ \vdots \\ \psi_d\end{bmatrix} = \psi_1 \ket{1} + \cdots + \psi_d\ket{d} = \sum_{i=1}^d \psi_i\ket{i}
\end{equation}
\item Bra:
\begin{equation}
\bra{\psi} = \ket{\psi}^\dagger = \ket{\psi^*}^T
\end{equation}
\item Braket:
\begin{equation}
\braket{\psi|\phi}
= \begin{bmatrix} \psi_1^* \cdots \psi_d^*\end{bmatrix} \begin{bmatrix} \phi_1\\ \vdots \\ \phi_d\end{bmatrix}
= \psi_1^*\phi_1 + \cdots + \psi_d^*\phi_d
\end{equation}
$\braket{\psi|\phi} $ is the hermitian product of $\psi$ and $\phi$.
\end{itemize}

\subsection{Measurement in a basis $B$}
$B$ is an orthonormal basis : $B = \{\ket{b_i}\}_{i=1}^d$ and $\forall i \braket{b_i|b_i} = \delta_{i, j}$.

\begin{figure}[h]
\centering
\begin{quantikz}
    \lstick{$\ket{\psi}$} & \qw & \meter{$B$} & \qw \arrow[r]
    & \rstick{$i$} \qw
\end{quantikz}
\caption{Circuit representation of the measurement of the state $\psi$}
\end{figure}

\noindent
The probability of the output of a measurement is given by the following formula :
\begin{equation}
    \mathbbm{P}(out=\ket{b_i}) = |\braket{b_i|\psi}|^2
\end{equation}
The physical object is projected into the state $\ket{b_i}$, the is physically called the "wave packet reduction".

\subsubsection{Qubit}
\begin{equation}
    \ket{0}= \begin{bmatrix} 1 \\ 0 \end{bmatrix}
    \hspace{1.2in}
    \ket{1}= \begin{bmatrix} 0 \\ 1 \end{bmatrix}
\end{equation}
\begin{equation}
    \ket{+} = \frac{1}{\sqrt2}(|0\rangle + |1\rangle)
    \hspace{0.7in}
    \ket{-}= \frac{1}{\sqrt2}(|0\rangle - |1\rangle)
\end{equation}

\subsubsection{Measurement in the basis $\{\ket{\pm}\}$}
\begin{figure}[h]
\centering
\begin{quantikz}
    \lstick{$\ket{0}$} & \qw & \meter{$\pm$} & \qw \arrow[r]
    & \rstick{$
        \begin{cases}
            +, & \text{w.p. $|\braket{+|0}|^2=\frac{1}{2}$}\\
            -, & \text{w.p. $|\braket{-|0}|^2=\frac{1}{2}$}
        \end{cases}$} \qw
\end{quantikz}
\caption{Measure of the state $\ket{0}$ in the basis $\ket{\pm}$}
\end{figure}

\subsection{Wiesner's Quantum Money}
Based on the conjugate coding.

\begin{itemize}
    \item \texttt{bills}:
    \begin{itemize}
        \item serial number
        \item a set of random qubit $E_r \in \{\ket{0}, \ket{1}, \ket{+}, \ket{-}\}^n$
        \item \texttt{mint} knows $\{\text{Serial Number}+\text{Random}\}$, sends it to the bank.
    \end{itemize}
    \item \texttt{Mint}: makes the bill, and gives it to the forger.
    \item \texttt{Forger}: tries to copy the bill, and spends the two to the bank.
    \item \texttt{Bank}: should accept the true one, reject the fake.
\end{itemize}

\begin{table}[h]
    \centering
    \begin{tabular}{c|c|c|c}
    \texttt{mint} & \texttt{forger} basis & \texttt{forger} m. & \texttt{bank} m. \\
    $\ket{0}$ & $\{\ket{0}, \ket{1}\}$ & $\ket{0}$ & $\ket{0} $\\
    $\ket{0}$ & $\{\ket{\pm}\}$ &
    $\begin{cases}
        \ket{+}, & \text{w.p. $\frac{1}{2}$}\\
        \ket{-}, & \text{w.p. $\frac{1}{2}$}
    \end{cases}$ &
    $\begin{cases}
        \ket{0}, & \text{w.p. $\frac{1}{2}$}\\
        \ket{1}, & \text{w.p. $\frac{1}{2}$}
    \end{cases}$
    \end{tabular}
\end{table}

We therefore deduce that
\begin{equation}
\mathbbm{P}(\text{get caugth})= 1-(1-\frac{1}{4})^n = 1-(\frac{1}{4})^n
\end{equation}

\subsection{Bennett and Brassard Quantum Key Exchange: BB84}
Goal: Alice and Bob $\rightarrow$ share a secret bit string , Eve does not know anything.
\\
Setting: Alice and Bob share a quantum channel and an authenticated classical channel.
\\\noindent
Steps:
\begin{enumerate}
    \item Alice prepares $n$ qubits $E_r \in \{\ket{0}, \ket{1}, \ket{+}, \ket{-}\}^n$, and she sends them to Bob
    \item Bob receives . He measure them in the basis $\{B_{0,1}, B_{+,-}\}$
    \item They use the public classical channel to compare the basis Bob used. They throw away the \textit{bad basis} qubits.
    \item Alice and Bob sample the data and compare the error rate $e$. If $e=0$, they keep the key;
        if $e = 25$, Eve kowns the key.
\end{enumerate}
What if $0 < e < 25$ ? Eve knows a part of the key.
\section{Unitary transformation}
A transformation is an isolated system, and it is reversible. \\
Let $T$ to be a transformation.
\begin{equation}
    \braket{T(\ket{\psi})|T(\ket{\psi})}=\braket{\psi|\psi}
\end{equation}
$T$ is linear.
\begin{equation}
    T(\alpha\ket{\psi} + \beta\ket{\phi}) = \alpha T(\ket{\psi}) + \beta T(\ket{\phi})
\end{equation}

$T$ acts like an unitary operator.
$T$ corresponds to a complex matrix $U$: $T(\ket{\phi}) = U\ket{\phi}$, $U \in \mathbbm{C}^{n \times n}$,
such that $U^\dagger U= Id$.

In the basis $\{\ket{i}\}_{i=0}^n$,$\braket{T(\ket{\psi})|T(\ket{\psi})} = \braket{i|j} = \delta_{i,j}$
%TODO : add the proof
\\
\\
We have :
\begin{itemize}[label=-]
    \item measurement in computational basis
    \item a machine making arbitrary unitary $U$
\end{itemize}
Let's build a measurement in basis $\{\ket{b_i}\}_i$

\begin{figure}[h]
    \centering
    \begin{quantikz}
        \lstick{$\ket{\psi}$} & \qw & \gate{U} & \phase{U\ket{\psi}}&\meter{} & \qw \arrow[r]
        & \rstick{$i$} \qw
    \end{quantikz}
    \caption{Circuit representation of the measurement unitary exptected behavior}
\end{figure}
\begin{equation}
    \mathbbm{P}(i) \eqdef |\braket{i|U|\psi}|^2 \eqgoal |\braket{b_i|\psi}|^2 \quad \forall \psi
\end{equation}

We want $\bra{i}U=\bra{b_i} \Leftrightarrow U^\dagger\ket{i} = \ket{b_i} \Leftrightarrow U = \sum_i \ket{i}\bra{b_i}$
% TODO : add the proof

Is $U$ an unitary ?

\begin{equation}
    \begin{aligned}
        U^\dagger U &= \Big(\sum_i \ket{b_i}\bra{i}\Big)\Big(\sum_j\ket{j}\bra{b_j}\Big) \\
        &= \sum_{i, j}\ket{b_i}\braket{i|j}\bra{b_j} \\
        &= \sum_i \ket{b_i}\bra{b_i} \\
        &= Id & U \text{ is an unitary. }
    \end{aligned}
\end{equation}

\section{Composition of systems}
Let $A \in \mathscr{H}_A = \mathbbm{C}^{d_A}$ and $B \in \mathscr{H}_B = \mathbbm{C}^{d_B}$
to be two systems in their respective vector spaces.
Then we can construct the space $\mathscr{H}_{AB} = \mathscr{H}_A \otimes \mathscr{H}_B$, its basis is
$\{\ket{ij}_{AB}\}_{i, j}$, and dim($\mathscr{H}_{AB}$) = $d_A \times d_B$
\\\noindent
If $\ket{\alpha} = \sum_i\alpha_i\ket{i}_A$ and $\ket{\beta} = \sum_i\beta_i\ket{i}_B$,
then $\ket{\phi}_{AB}=\ket{\alpha} \otimes \ket{\beta} = \sum_{i,j}\alpha_i\beta_j\ket{i}_A\ket{j}_B$,
and $\ket{\phi}_{AB} \in \mathscr{H}_{AB}$.
\\\noindent
$\ket{\phi}_{AB}$ is a joint state of systems $A$ and $B$.
\\ \\\noindent
Not all states of $\mathscr{H}_{AB}$ are separables into one state of $\mathscr{H}_{A}$ and one state of $\mathscr{H}_{B}$

For example : $\frac{1}{\sqrt{2}}(\ket{00} + \ket{11}) \ \mathscr{H}_{AB}$, but
$\nexists \ket{\alpha}\in \mathscr{H}_A, \ket{\beta}\in \mathscr{H}_B$, such that \\\noindent
$\ket{\alpha}\otimes\ket{\beta} = \frac{1}{\sqrt{2}}(\ket{00} + \ket{11})$.

%TODO : add the part about the partial measurement
\subsection{No cloning theorem}
\subsubsection*{The no cloning theorem}
\begin{equation}
    \text{There is no } U \text{ such that } \forall \ket{\psi} \in \mathscr{H}, U\ket{\psi} = \ket{\psi} \otimes \ket{\psi}
    \in \mathscr{H} \otimes \mathscr{H}.
\end{equation}
\subsubsection*{Proof}
Suppose there exists a such unitary U, then U in a cloning operator.
\begin{equation}
    \begin{aligned}
        U\ket{0} \eqdef \ket{0}\ket{0} \\
        U\ket{1} \eqdef \ket{1}\ket{1}
    \end{aligned}
\end{equation}

By computing the application of $U$ on the state $\ket{+}$, we get on the one hand, by linearity of unitaries
\begin{equation}
    U(\frac{\ket{0}+\ket{1}}{\sqrt{2}}) = \frac{\ket{00}+\ket{11}}{\sqrt{2}}
\end{equation}

and on the other hand, by definition of the operator behavior
\begin{equation}
    U(\frac{\ket{0}+\ket{1}}{\sqrt{2}}) = \frac{\ket{0}+\ket{1}}{\sqrt{2}} \otimes \frac{\ket{0}+\ket{1}}{\sqrt{2}}
    =\frac{1}{2}(\ket{00}+\ket{01}+\ket{10}+\ket{11})
\end{equation}

which is a contradiction. Then such a $U$ operator can not exist.
\section{Measurements}
\subsection{Projective measurement}
A projective measurement is described by an observable, a Hermitian operator. They are defined by a set of projectors $\{\Pi_j\}_{j=1}^k, k \leq d$.\\
Projectors properties:
\begin{equation}
    \forall j, \Pi_j^2 = \Pi_j \qquad \Pi_j\Pi_i = \delta_{i,j}\Pi_j
\end{equation}
\\\noindent
A projector is defined as follows:
\begin{equation}
    \Pi_j = \sum_{l=1}^{d_j}\ket{l_l^j}\bra{l_l^j}
\end{equation}
Upon measuring the state $\ket{\psi}$, the probability of getting result $j$ us given by
\begin{equation}
    \braket{\psi|\Pi_j|\psi} = \| \Pi_j\ket{\psi}\|^2
\end{equation}
Given that outcome $j$ occured, the state of the quantum system immediately after the measurement is
\begin{equation}
    \frac{\Pi_j\ket{\psi}}{\| \Pi_j\ket{\psi}\|^2}
\end{equation}

\subsection{Observables}
Observables correspond to physical quantities, with values in $\mathbbm{R}$.
They are well defined in a basis $\{\ket{b_i}\}_i$ (i.e $\forall \ket{b_i}, \exists a_i \in \mathbbm{R})$

\paragraph{Note :} $\alpha\ket{b_1} + \beta\ket{b_2}$ has \textbf{not always} a well defined value.

An observable is defined as follow:
\begin{equation}
    O \eqdef \sum_i o_i \underbrace{\ket{b_i}\bra{b_i}}_{\text{projector on $\ket{b_i}$}}
    = \sum_j o_j \Pi_j
\end{equation}
$O$ is diagonalizable by definition and $O^\dagger = O$: $O$ is hermitian.
$$
\text{Shape of $O$ } :
\begin{pmatrix}
    o_1 & 0 & \cdots & 0 \\
    0 & o_2 & \cdots & \vdots \\
    \vdots  & \vdots  & \ddots & \vdots  \\
    0 & \cdots & \cdots & a_d
\end{pmatrix}
$$

The probability of getting the result $i$ by measuring $O$ on a state $\ket{\psi}$ is $\braket{\psi|\Pi_i|\psi}$
\subsubsection*{Average value of an observable}

The average value of t of $O$, written $\langle O\rangle$, is by definition given by
\begin{equation}
    \begin{aligned}
        \langle O\rangle
            & = \sum_i o_i \mathbbm{P}(i\ket{\psi}) \\
            & = \sum_i o_i \|\Pi_i\ket{\psi}\|^2 \\
            & = \sum_i o_i \braket{\psi|\Pi_i|\psi} \\
            & = \bra{\psi}\sum_i o_i \Pi_i \ket{\psi} \\
            & = \braket{\psi | O |\psi}
    \end{aligned}
\end{equation}

From this formula for the averag follows a formula for the standard deviation associated to the observation of $O$
\begin{equation}
    \Delta^2 O = \langle(O - \langle O \rangle)^2\rangle = \langle O^2 \rangle - \langle O \rangle^2
\end{equation}

\subsubsection*{Example}
Using the Pauli matrix $\sigma_x = X =
    \begin{bmatrix}
    0 & 1 \\
    1 & 0
    \end{bmatrix}
    = \ket{+}\bra{+} - \ket{-}\bra{-}
    $.
\\\noindent
Known results : $X\ket{+} = \ket{+}$ and $X\ket{-}=-\ket{-}$.
\\\noindent
We define $\ket{\theta}:=\cos\omega\ket{0}+\sin\omega\ket{1}$
\\\noindent
Then
\begin{equation}
    \begin{aligned}
        \langle X \rangle_{\ket{\theta}}
            & = \braket{\theta | X |\theta} \\
            & = [\cos \theta \sin \theta]
                \begin{bmatrix}0 & 1 \\ 1 & 0 \end{bmatrix}
                \begin{bmatrix} \cos \theta\\ sin \theta \end{bmatrix} \\
            & = 2 \sin\theta\cos\theta \\
            & = sin 2\theta
    \end{aligned}
\end{equation}


\subsection{Generalized measurements}
A generalized measurement is defined by
\begin{equation}
    \{K_i\}_i \text{ such that } \sum_i K_i^\dagger K_i=Id
\end{equation}
where the $K_i$ are called Kraus Operators.
The probability of getting the result $i$ from a general measurement operator is given by
$\mathbbm{P}(i) = \|K_i\ket{\psi}\|^2$
, and the state of the system just after the measurement is $K_i\ket{\psi}=\frac{K_i\ket{\psi}}{\|K_i\ket{\psi}\|^2}$

\subsubsection*{Generalized measurement $\rightarrow$ Operator}
If $i \in \{1\}$ then
$K_1^\dagger K_1=Id \Rightarrow K_1$ is unitary.%FIXME: not sure of that.

\subsubsection*{Generalized measurement $\rightarrow$ Set of projectors}
If $K_i := \Pi_i$ then $\sum_i K_i^\dagger K_i = \sum_i \Pi_i^\dagger\Pi_i = \sum_i \Pi_i = Id$

\subsubsection*{Example}
With prob. $P_j$, I measure $\{\Pi_{ij}\}_i \text{ } (\sum_i \Pi_{ij} = Id)$ and I measure $U_j$ on the output state.
%TODO : add the scheme
Probability of getting $ij$ :
\begin{equation}
    \begin{aligned}
        \mathbbm{P}(ij)
            &= P_j\braket{\psi | \Pi_{ij} U^\dagger U \Pi_{ij} | \psi} \\
            &= P_j\braket{\psi|\Pi_{ij}|\psi}
    \end{aligned}
\end{equation}
\noindent
And the resulting state is $\dfrac{U\Pi_{ij}\ket{\psi}}{\|\Pi_{ij}\ket{\psi}\|^2}$

\noindent
Let $\{K_{ij} =\sqrt{P_j} U \Pi_{ij}\}_{ij}$, then
\begin{equation}
    \begin{aligned}
        \sum_{ij} K_{ij}^\dagger K_{ij}
            & = \sum_{ij} P_j \Pi_{ij} U^\dagger U \Pi_{ij} \\
            & = \sum_j P_j \sum_i \Pi_{ij} \\
            & = \sum_j P_j \\
            & = Id
    \end{aligned}
\end{equation}
\noindent
Can we associate each set $\{K_i\}_i$ with a $U$ and a $\{\Pi_i\}_i$ ?
\begin{figure}[h]
    \centering
\begin{quantikz}
        \lstick{$\ket{\psi}_I$}
        & \qw \gategroup[2,steps=5,style={dashed,rounded corners,fill=blue!20, inner xsep=2pt},background]{{\sc $\{K_i\}_i$}}
        & \qw
        & \qw
        & \gate[2]{U}
        & \qw\arrow[r]
        &
        & \text{ouput $o$}
        \\
        &
        &
        & \lstick{   $\ket{0}_A$}
        & \qw
        & \meter{M$\{\ket{i}\}$}
        & \qw \arrow[r]
        & \text{ }i
    \end{quantikz}
    \caption{Circuit representation of such $U$ and $\{\Pi_i\}_i$. \\Note : $\mathscr{H}_A\otimes \mathscr{H}_I = \mathscr{H}_O\otimes\mathscr{H}_M$}
\end{figure}
\\\noindent
$\forall i,$ the output state of the system is
\begin{equation}
    \label{output-state-kraus}
    (I_o \otimes \ket{i}_M\bra{i})U\ket{\psi}\otimes\ket{0}_A = \ket{i}_M\braket{i|U|0}_A\ket{\psi}_I
\end{equation}
Assume $K_i = _M\braket{i|U|0}_A$. With (\ref{output-state-kraus}), we deduce that the output state is
$K_i\ket{\psi}$, w.p. $\braket{\psi|K_i^\dagger K_i|\psi}$.
Is that a valid set of operators $\{K_i\}_i$ ?

\begin{equation}
    \begin{aligned}
        \sum_i K_i^\dagger K_i
        & = \sum_i (_A\bra{0}\otimes I_I) U^\dagger (\ket{i}_M \otimes I_O)(I_O \otimes _M\bra{i})U(I_I\otimes \ket{0}_A)\\
        & = (_A\bra{0}\otimes I_I)
        \underbrace{U^\dagger (
        \underbrace{
            \underbrace{
                \sum_i \ket{i}_M}_{=I_M}
                    \otimes I_O
            }_{=I_{OM}})
            (I_O \otimes _M\bra{i})U
        }_{= I_{OA}}(I_I\otimes \ket{0}_A) \\
        & = (_A\bra{0}\otimes I_I)I_{OA}(I_I\otimes \ket{0}_A) \\
        & = I_O \qquad \text{$K_i$ is a valid set.}
    \end{aligned}
\end{equation}

\subsubsection*{$\{K_i\}_i \rightarrow$ Unitary}
Let $U := \sum_i K_i \otimes \ket{i}_{MA}\bra{0} + \cdots $.
The $\cdots$ represent extra terms used to make $U$ a unitary, but can be neglected in the computation.
By tensoring with $\ket{0}_A$, we obtain
\begin{equation}
    U\ket{\psi}\otimes \ket{0}_A = \sum_i K_i \ket{\psi}\otimes\ket{i}
\end{equation}
And then
\begin{equation}
    \begin{aligned}
        _A\braket{0 |U^\dagger U|0}_A
            & = _A\bra{0}(\sum_i\ket{0}_{AM}\bra{i} K_i^\dagger \cdot \sum_j K_j\ket{j}_{AM}\bra{0})\ket{0}_A \\
            & = \underbrace{_A\braket{0|0}_A}_{=1} \cdot \sum_{ij}(_M\bra{i}\otimes K_i^\dagger)(\ket{j}_M\otimes K_j) \underbrace{_A\braket{0|0}_A}_{=1} \\
            & = \sum_{ij}\underbrace{\braket{i|j}}_{\delta_{ij}} K_i^\dagger K_j \\
            & = \sum_i K_i^\dagger K_i \\
            & = Id
    \end{aligned}
\end{equation}
\subsection{POVMs}
POVMs means Projective Operator Valued Measure : we "only care about the output $i$, not the state $K_i\ket{\psi}$".
The probability of getting $i$ is $\braket{\psi | K_i^\dagger K_i | \psi}$.\\
Let $E_i = K_i^\dagger K_i$.
POVMs are then defined by the set $\{E_i\}_i$, such that $\sum_i E_i = Id, E_i\geq 0$.\\

$E_i$ is semi-definite positive: $ \forall \psi, \braket{\psi|E_i|\psi} \geq 0$. This implies that $E_i$ is hermitian, and all its eigenvalues are $\geq 0$.
\subsubsection*{POVM $\rightarrow$ Kraus operators}
Let $K_i = \sqrt{E_i}$. $\{K_i\}_i$ is a well defined set of operators.
\end{document}